% Laget automatisk av elevark/lag_elevark.py fra sim/usikkerhet-beregning.html. Endringer her blir overskrevet.
\documentclass[11pt]{article}
\usepackage{elevark}
\begin{document}
\elevarkhode{Fysikk 2 · Tenke og gjøre fysikk}{Usikkerhet i beregnede størrelser}{https://fysikklab.no/sim/usikkerhet-beregning.html}{Simuleringen}

Når vi regner ut $g$ fra lengden og svingetiden til en pendel, har begge målingene usikkerhet. Hvor stor blir usikkerheten i $g$? Og hvilken måling bør vi gjøre bedre for å få et sikrere resultat?

\section*{Oppgaver}

\begin{oppgaveliste}
\oppgave \fase{test}{Regn}En pendel er $L = (0{,}800 \pm  0{,}005)\,\mathrm{m}$ lang og har svingetid $T = (1{,}79 \pm  0{,}02)\,\mathrm{s}$. Regn ut $g$ og den relative usikkerheten i $g$ med den enkle regelen. Sjekk.
\svar
\oppgave \fase{forklar}{Forklar}Hvorfor teller usikkerheten i $T$ dobbelt? Hvilken måling bidrar mest til usikkerheten i $g$?
\svar
\oppgave \fase{test}{Undersøk}Eleven måler i stedet tiden for 10 svingninger og får 17,90~s ± 0,02~s. Hva blir usikkerheten i $T$ nå, og hva blir usikkerheten i $g$? Skriv det inn og sjekk.
\svar
\oppgave \fase{test}{Undersøk}Trykk på «Simuler 5000 målinger» for pendelen. Ligger spredningen nærmest den enkle regelen eller kvadratsummen? Hvorfor gir den enkle regelen et større tall?
\svar
\oppgave \fase{test}{Regn}Velg tetthet til en kloss. Hvilken måling bidrar mest? Hvorfor er det lurt å måle den minste lengden med et skyvelære i stedet for en linjal?
\svar
\oppgave \fase{forklar}{Forklar}Når to tall legges sammen eller trekkes fra, legger vi sammen de absolutte usikkerhetene. Vis med et eksempel hvorfor det kan gi stor relativ usikkerhet å trekke to like store tall fra hverandre.
\svar
\oppgave \fase{utfordring}{Utfordring}I forsøket med fritt fall er falltiden 0,553~s. Hvor nøyaktig må tiden måles for at $g$ skal få under 1~\% usikkerhet? Hvilket måleutstyr trenger dere?
\svar
\end{oppgaveliste}

\end{document}
